% ============================================================
% Section 2: Pre-Classical Mathematics (Slides 4-14)
% ~20,000 BCE -- 500 BCE
% Ishango bone, Babylon, Egypt, Maya
% ============================================================
\section{Pre-Classical Mathematics}

% ----------------------------------------------------------
% Slide 4: Section Divider
% ----------------------------------------------------------
{
\setbeamercolor{background canvas}{bg=babylonBrown}
\begin{frame}[plain,shrink=16]
  \centering
  \vspace{1cm}
  \textcolor{white}{\Huge\bfseries Where It All Began}\\[1em]
  \textcolor{white!80}{\large Pre-Classical Mathematics\quad $\sim$20,000 BCE -- 500 BCE}\\[1.5em]

  % Mini-map: Africa, Fertile Crescent, Nile, Mesoamerica
  \visualplaceholder[5cm]{Region highlights}

  \vspace{0.5em}
  % Timeline marker
  \visualplaceholder[5cm]{Diagram}

  \note{
    Section transition. Emphasise geographic spread: mathematics did not start
    in one place. Point to Central Africa (Ishango), Mesopotamia, the Nile, and Mesoamerica.
  }
\end{frame}
}

% ----------------------------------------------------------
% Slide 5: The Ishango Bone
% ----------------------------------------------------------
\begin{frame}[shrink=13]{The Ishango Bone --- The Oldest Mathematics?}

  \begin{columns}[T]
    \begin{column}{0.55\textwidth}
      \begin{personbox}[Ishango Bone]{babylonBrown}
        \small
        \textbf{Date:} c.\ 20,000 BCE\\
        \textbf{Discovered:} 1960, near Lake Edward, modern DR Congo\\
        \textbf{Now held at:} Royal Belgian Institute of Natural Sciences, Brussels
      \end{personbox}

      \vspace{0.5em}
      \begin{itemize}
        \item A bone tool with \alert{notched markings}
        \item Possibly represents: arithmetic sequences, prime numbers, or a \textbf{lunar calendar}
        \item One of the \textbf{earliest known mathematical artefacts}
        \item Interpretation is \emph{debated} (De Heinzelin, Marshack vs.\ sceptics)
      \end{itemize}
    \end{column}

    \begin{column}{0.42\textwidth}
      \visualplaceholder{Photograph of the Ishango bone showing three columns of tally marks (Royal Belgian Institute of Natural Sciences)}

      \vspace{0.3em}
      % Tally mark schematic
      \visualplaceholder[5cm]{Complex diagram}
    \end{column}
  \end{columns}

  \note{
    ``Mathematics may be as old as humanity itself. 20,000 years before Babylon,
    someone in Central Africa was scratching mathematical patterns into bone.''
    \verifytag{Interpretation is debated. Some scholars (Marshack, De~Heinzelin) see
    mathematical significance; others see it as a simple tally or tool handle.
    Present as ``possibly'' mathematical.}
  }
\end{frame}

% ----------------------------------------------------------
% Slide 6: Babylon -- The First Mathematicians
% ----------------------------------------------------------
\begin{frame}[shrink=11]{Babylon --- The First Mathematicians}

  \begin{columns}[T]
    \begin{column}{0.52\textwidth}
      \begin{itemize}
        \item Civilization in Mesopotamia (modern Iraq), c.\ 2000--500 BCE
        \item Developed a \alert{base-60} (sexagesimal) number system
        \item Why \textbf{60 minutes} in an hour, \textbf{360 degrees} in a circle
        \item Wrote mathematics on \textbf{clay tablets} in cuneiform script
        \item Thousands of mathematical tablets survive
      \end{itemize}

      \vspace{0.5em}
      \begin{insightbox}
        The next time you check the clock, thank a Babylonian scribe from 4,000 years ago.
      \end{insightbox}
    \end{column}

    \begin{column}{0.45\textwidth}
      \visualplaceholder{Photograph of a cuneiform clay tablet with mathematical content (Yale Babylonian Collection or British Museum)}

      \vspace{0.3em}
      % Base-60 numeral demonstration
      \visualplaceholder[5cm]{Wedge for 1}
    \end{column}
  \end{columns}

  \note{
    Establish Babylon as the cradle of systematic mathematics.
    Emphasise the base-60 system and its legacy in our daily lives.
    Transition: ``And they did more than just count---they solved equations.''
  }
\end{frame}

% ----------------------------------------------------------
% Slide 7: Plimpton 322
% ----------------------------------------------------------
\begin{frame}[shrink=25]{Plimpton 322 --- A 3,800-Year-Old Mystery}

  \begin{columns}[T]
    \begin{column}{0.55\textwidth}
      \begin{personbox}[Anonymous Babylonian Scribe(s)]{babylonBrown}
        \small
        \textbf{Date:} c.\ 1800 BCE\\
        \textbf{Origin:} Mesopotamia (modern Iraq)\\
        \textbf{Now held at:} Columbia University Rare Book \& Manuscript Library
      \end{personbox}

      \vspace{0.4em}
      \begin{itemize}
        \item Contains \alert{Pythagorean triples}
        \item Written $\sim$\textbf{1,200 years BEFORE Pythagoras}
        \item Interpretation debated:
          \begin{itemize}
            \item Neugebauer \& Sachs (1945): table of triples
            \item Robson (2001): teacher's aid for exercises
          \end{itemize}
      \end{itemize}

      \vspace{0.3em}
      \begin{insightbox}
        The theorem we call ``Pythagorean'' was known to Babylonians over a millennium earlier. Math rarely has a single inventor.
      \end{insightbox}
    \end{column}

    \begin{column}{0.42\textwidth}
      \visualplaceholder{High-resolution photo of Plimpton 322 tablet (Columbia University, public domain)}

      \vspace{0.3em}
      % Small table of triples
      {\small
      \begin{tabular}{ccc}
        \toprule
        $a$ & $b$ & $c$ \\
        \midrule
        119 & 120 & 169 \\
        3367 & 3456 & 4825 \\
        65 & 72 & 97 \\
        \bottomrule
      \end{tabular}
      }\\[0.2em]
      {\tiny\textcolor{gray}{Selected rows from Plimpton 322}}
    \end{column}
  \end{columns}

  \note{
    Emphasise that ``Pythagorean'' is a misnomer in terms of priority.
    \verifytag{Exact interpretation of Plimpton 322 is debated.
    Neugebauer \& Sachs (1945) = Pythagorean triples;
    Robson (2001) = teacher's aid. Both should be mentioned.}
    ``Math rarely has a single inventor.''
  }
\end{frame}

% ----------------------------------------------------------
% Slide 8: Egyptian Mathematics -- Practical Genius
% ----------------------------------------------------------
\begin{frame}[shrink=5]{Egyptian Mathematics --- Practical Genius}

  \begin{columns}[T]
    \begin{column}{0.52\textwidth}
      \begin{itemize}
        \item Driven by \textbf{practical needs}:
          \begin{itemize}
            \item Land surveying after Nile floods
            \item Architecture (pyramids!)
            \item Taxation and accounting
          \end{itemize}
        \item Two main sources:
          \begin{itemize}
            \item \alert{Rhind Papyrus} (next slide) --- larger, more complete
            \item \alert{Moscow Papyrus} --- contains the frustum formula
          \end{itemize}
        \item Used unit fractions: $\frac{2}{5} = \frac{1}{3} + \frac{1}{15}$
      \end{itemize}
    \end{column}

    \begin{column}{0.45\textwidth}
      \visualplaceholder{Section of the Rhind Papyrus (British Museum, EA 10057-8)}

      \vspace{0.5em}
      % Pyramid sketch
      \visualplaceholder[5cm]{height}
    \end{column}
  \end{columns}

  \note{
    Transition note: The Rhind Papyrus and the Moscow Papyrus are our two main sources.
    We will meet the scribe Ahmes on the next slide, then look at a specific geometry problem.
    Egyptian math was supremely practical --- no ``proofs'' in the Greek sense, but brilliantly effective.
  }
\end{frame}

% ----------------------------------------------------------
% Slide 9: Ahmes (Ahmose)
% ----------------------------------------------------------
\begin{frame}[shrink=4]{Ahmes --- The Oldest Named Mathematician}

  \begin{columns}[T]
    \begin{column}{0.55\textwidth}
      \begin{personbox}[Ahmes (Ahmose)]{babylonBrown}
        \small
        \textbf{Dates:} fl.\ c.\ 1550 BCE\\
        \textbf{Origin:} Ancient Egypt\\
        \textbf{Contribution:} Scribe who copied (and possibly compiled) the
        \emph{Rhind Mathematical Papyrus} --- our best source on Egyptian mathematics
      \end{personbox}

      \vspace{0.4em}
      \begin{itemize}
        \item Contains problems in \textbf{arithmetic}, \textbf{algebra}, \textbf{geometry}
        \item Implies $\pi \approx 3.16$ (via circle area formula)
        \item Ahmes wrote: \emph{``Accurate reckoning for inquiring into things, and the knowledge of all things, mysteries\ldots\ all secrets.''}
      \end{itemize}
    \end{column}

    \begin{column}{0.42\textwidth}
      \visualplaceholder{Portion of the Rhind Papyrus showing a geometry problem (British Museum digital collection)}
    \end{column}
  \end{columns}

  \note{
    Ahmes is the oldest person we know by name who did mathematics.
    He was a scribe, not necessarily the inventor of the methods.
    \verifytag{fl.\ c.\ 1550 BCE is the standard dating. The ``accurate reckoning''
    quote translation should be double-checked against the British Museum source.}
    ``He is copying and perhaps editing work that was already old in his time.''
  }
\end{frame}

% ----------------------------------------------------------
% Slide 10: Egyptian Geometry in Action [EXTENDED]
% ----------------------------------------------------------
\begin{frame}[shrink=16]{Egyptian Geometry in Action}

  \begin{columns}[T]
    \begin{column}{0.48\textwidth}
      \textbf{Circle area (Rhind Papyrus):}\\[0.3em]
      % Octagon approximation diagram
      \visualplaceholder[5cm]{Square 9x9 conceptual (scaled to 4.5)}

      \vspace{0.3em}
      {\small Egyptian rule: $A \approx \left(\frac{8}{9}d\right)^2$\\
      $\Rightarrow \pi \approx \frac{256}{81} \approx 3.16$}
    \end{column}

    \begin{column}{0.48\textwidth}
      \textbf{Frustum volume (Moscow Papyrus):}\\[0.3em]
      % Frustum diagram
      \visualplaceholder[5cm]{Frustum (trapezoid cross-section)}

      \vspace{0.3em}
      {\small $V = \frac{h}{3}(a^2 + ab + b^2)$}\\[0.3em]
      {\small \textcolor{accentOrange}{Exactly correct!}\\
      c.\ 1850 BCE --- from the Moscow Papyrus}

      \vspace{0.5em}
      \begin{insightbox}
        They got $\pi$ accurate to within 1\% --- without calculus, without computers, using reeds and papyrus.
      \end{insightbox}
    \end{column}
  \end{columns}

  \note{
    Two remarkable results. The octagon approximation for circle area gives
    $\pi \approx 3.16$, which is within 1\% of the true value.
    The frustum formula is \emph{exact} --- we still use it today.
    ``How did they derive an exact formula without algebraic notation or proof?
    We don't know --- the papyrus just states it as a recipe.''
  }
\end{frame}

% ----------------------------------------------------------
% Slide 11: Mesopotamian Algebra [EXTENDED]
% ----------------------------------------------------------
\begin{frame}[shrink=11]{Mesopotamian Algebra}

  \begin{columns}[T]
    \begin{column}{0.50\textwidth}
      \textbf{Babylonians solved quadratic equations}\\[0.3em]
      {\small (in verbal/geometric form, not symbolic notation)}

      \vspace{0.6em}
      \textbf{Babylonian word problem:}\\
      {\small\itshape ``I have added the area and the side of my square: 0;45.''}

      \vspace{0.5em}
      \textbf{Modern equivalent:}\\
      $x^2 + x = \frac{3}{4}$

      \vspace{0.5em}
      {\small Babylonian solution (completing the square):}\\
      $x = \sqrt{\frac{3}{4} + \frac{1}{4}} - \frac{1}{2} = \frac{1}{2}$
    \end{column}

    \begin{column}{0.47\textwidth}
      % Geometric completing the square
      \visualplaceholder[5cm]{Adding the side: rectangle x*1}

      \vspace{0.5em}
      \begin{insightbox}
        Your algebra homework has a 4,000-year pedigree.
      \end{insightbox}
    \end{column}
  \end{columns}

  \vspace{0.3em}
  \begin{center}
    \colorbox{accentOrange!15}{\parbox{0.85\textwidth}{\centering\small
      \faUsers\ \textbf{Think--Pair--Share:} ``Why do you think the Babylonians never developed
      the concept of mathematical proof? Was their approach inferior, or just different?'' \emph{(2--3 min)}
    }}
  \end{center}

  \note{
    Interactive moment! After presenting the algebra, pause for Think--Pair--Share.
    Give students 1 minute to think, 1 minute to discuss with a neighbour,
    then take 2--3 responses.
    Good answers might include: ``they didn't need proof because the recipes worked''
    or ``their math was for practical use, not abstract truth.''
    Transition to the next slide which addresses this directly.
  }
\end{frame}

% ----------------------------------------------------------
% Slide 12: What Was Missing?
% ----------------------------------------------------------
\begin{frame}{What Was Missing?}

  \begin{center}
    {\large Ancient mathematics was \alert{computational}, not \alert{axiomatic}.}
  \end{center}

  \vspace{0.5em}
  \begin{columns}[T]
    \begin{column}{0.47\textwidth}
      \visualplaceholder[5cm]{Diagram}
    \end{column}

    \begin{column}{0.06\textwidth}
      \centering\vspace{1cm}
      {\Large\textcolor{gray}{vs.}}
    \end{column}

    \begin{column}{0.47\textwidth}
      \visualplaceholder[5cm]{Diagram}
    \end{column}
  \end{columns}

  \vspace{0.8em}
  \begin{center}
    {\small This is \textbf{not a deficiency} --- the procedural approach served its purpose
    brilliantly for millennia.\\
    But the Greeks would add something new: the idea that you must \emph{prove} your claims.}
  \end{center}

  \note{
    Key conceptual slide. The point is NOT that the Greeks were ``better'' but that
    they introduced a fundamentally different way of thinking about mathematics.
    Both approaches have value. Modern mathematics uses both: we compute AND we prove.
  }
\end{frame}

% ----------------------------------------------------------
% Slide 13: Mayan Mathematics
% ----------------------------------------------------------
\begin{frame}[shrink=25]{Mayan Mathematics --- An Independent Tradition}

  \begin{columns}[T]
    \begin{column}{0.52\textwidth}
      \begin{personbox}[Maya Civilization]{babylonBrown}
        \small
        \textbf{Era:} c.\ 4th century CE and earlier\\
        \textbf{Region:} Mesoamerica (modern Mexico, Guatemala, Belize, Honduras)
      \end{personbox}

      \vspace{0.4em}
      \begin{itemize}
        \item \alert{Vigesimal} (base-20) number system
        \item Independently invented \textbf{the concept of zero}
        \item Long Count calendar required sophisticated positional arithmetic
        \item Zero used \textbf{centuries before Europe} knew the concept
      \end{itemize}

      \vspace{0.3em}
      \begin{insightbox}
        Zero was invented independently at least twice: in India and in Mesoamerica.
      \end{insightbox}
    \end{column}

    \begin{column}{0.45\textwidth}
      % Mayan numeral system
      \visualplaceholder[5cm]{Mayan Numerals}

      \vspace{0.3em}
      \visualplaceholder[4.5cm]{Mayan zero glyph and section of a Long Count calendar stone (Wikimedia Commons)}
    \end{column}
  \end{columns}

  \note{
    ``Zero was invented independently at least twice in human history --- in India and
    in Mesoamerica. The Maya used it centuries before Europe even knew the concept existed.''
    \verifytag{Exact dating of the earliest Mayan zero is debated. Tres Zapotes stele
    (Epi-Olmec, 36 BCE) may contain an early Long Count date; the classic Mayan zero
    glyph appears firmly by the 4th century CE.}
  }
\end{frame}

% ----------------------------------------------------------
% Slide 14: Legacy of Ancient Mathematics
% ----------------------------------------------------------
\begin{frame}{Legacy of Ancient Mathematics}

  \begin{columns}[T]
    \begin{column}{0.50\textwidth}
      \textbf{What survives from the ancient world:}
      \begin{itemize}
        \item \textbf{Base-60 system} \textrightarrow\ clocks, angles
        \item \textbf{Basic algebra} \textrightarrow\ quadratic equations
        \item \textbf{Geometry} \textrightarrow\ surveying, architecture
        \item \textbf{Approximations of $\pi$}
        \item \textbf{Independent number systems} across continents
        \item \textbf{Zero} appearing in multiple civilizations
      \end{itemize}
    \end{column}

    \begin{column}{0.47\textwidth}
      % Infographic: ancient math in daily life
      \visualplaceholder[5cm]{Clock icon}
    \end{column}
  \end{columns}

  \vspace{0.5em}
  \begin{center}
    \textcolor{greekBlue}{\emph{``But one civilization would transform math from a tool
    into a way of thinking\ldots''}}
  \end{center}

  \note{
    Summary slide. Emphasise that ancient mathematics was GLOBAL, not just Babylonian/Egyptian.
    The Ishango bone, Mayan zero, and Egyptian geometry all arose independently.
    Transition: ``But one civilisation would transform math from a tool into a way of thinking\ldots''
    This leads to the Greek section.
  }
\end{frame}
